\begin{pastexam}[label={px:ncku-109-prog-04}]{成大資訊聯招 109 學年度 演算法第 1 題}
（問答）

Is $5^{n+1}=O(5^n)$? Is $6^{2n}=O(6^n)$? (10\%)
\end{pastexam}

\begin{pastsoln}
答案：$5^{n+1} = O(5^n)$ 成立；$6^{2n} = O(6^n)$ 不成立。

照定義（資結書第~1.5~節）：$f(n) = O(g(n))$ 若且唯若存在正的常數 $c$、$n_0$，使得所有
$n \ge n_0$ 都有 $0 \le f(n) \le c\,g(n)$。

\textbf{$5^{n+1}$}。 $5^{n+1} = 5 \cdot 5^n$。 取 $c = 5$、$n_0 = 1$，所有 $n \ge 1$ 都有
$5^{n+1} = 5 \cdot 5^n \le 5 \cdot 5^n$，所以 $5^{n+1} = O(5^n)$。 指數加上一個常數，只是
乘上一個常數倍。

\textbf{$6^{2n}$}。 $6^{2n} = (6^2)^n = 36^n$。 假設有常數 $c$、$n_0$ 使所有 $n \ge n_0$ 都有
$36^n \le c \cdot 6^n$，兩邊除以 $6^n$ 得 $6^n \le c$。 但 $6^n$ 隨 $n$ 無限增大，取
$n \ge \max(n_0, \log_6 c + 1)$ 就有 $6^n > c$，矛盾。 所以 $6^{2n} \ne O(6^n)$。 指數乘上
常數倍，底數變成原來的次方，是不同的成長級數。

兩題的差別可以用比值看：$5^{n+1}/5^n = 5$ 是常數；$6^{2n}/6^n = 6^n \to \infty$。
\end{pastsoln}
